\section{Definition of the Study Context}
\subsection{Country Selection for Comparative Analysis}
This study evaluates four European countries: France, Italy, Poland and Switzerland, to assess how national energy contexts influence the ranking of nuclear reactor designs.

\paragraph{Italy}'s nuclear landscape remains constrained by the 1987 post-Chernobyl referendum which mandated complete phase-out and its 2011 reaffirmation following Fukushima. However, since Italian referendums are abrogative rather than prescriptive, government actions are enough to reverse these policies, on this premis, the current administration seems willingly to reintroduce nuclear power \cite{Sole24Ore_Italy_Nuclear_Strategy_2024}. Italy's history of political volatility and the long-term nature of nuclear projects create significant implementation challenges regardless of public opinion trends. 

For this study we have taken into consideration only the northen part of Italy due to the country's fragmented electricity market structure. The north bidding zone is the most energy intensive with an average load of 17.66 $\pm$ 4.35 GW and a sigificant dependence on imports, which averaged 5.23 $\pm$ 1.67 GW in 2024. Imports therfore account for almost 30\% of its electric demand and cross-border flows show that imports are active 99\% hours of the year. 

\paragraph{Poland} shows strong public backing for nuclear energy development, with surveys reporting up to $92\%$ support \cite{NucNet_Poland_Nuclear_Support_2024}. The public sentiment aligns with tangible progress in the country's nuclear program, including preliminary site works and vendor agreements \cite{Westinghouse_Poland_Nuclear_2025,WorldNuclearNews_Poland_Preliminary_Works_2025}, suggesting robust governmental commitment to nuclear deployment. The Polish grid produced 55\% of its electricity with coal \cite{IEA_Electricity_2024}, making it the country with the highest green house gas emission intenisty in Europe \cite{ElectricityMaps_CarbonIntensity_2024}.


\paragraph{Switzerland} presents a more complex picture. Following the 2011 Fukushima accident, the government adopted a phase-out policy later confirmed by a 2017 referendum ($42\%$ approval) \cite{WorldNuclear_Switzerland_Profile_2025}. However, recent surveys indicate shifting public opinion, with $53\%$ now supporting nuclear power \cite{Euronuclear_Swiss_Nuclear_Poll_2023,NucNet_Nuclear_Support_Survey_2023,WorldNuclearNews_Swiss_Nuclear_Poll_2023}. The federal government has subsequently signaled potential reconsideration of nuclear's role in meeting climate targets \cite{swissinfo_Nuclear_Power_Switzerland_2025}, though the narrow margin of support and Switzerland's direct democracy system create policy uncertainty. 

The Swiss grid is characterized by a low demand which averaged 6.79 $\pm$ 1.12 GW in 2024, and generation dominated by renewables, 65\% average share, particularly hydroelectric which provided 59\% of the electricity in 2024 \cite{IEA_Electricity_2024}. This high dependence on renewables is reflected on substantial cross border flow variability which average 2.71 $\pm$ 3.14 GW. The large standard deviation reflects periodic exports due to overproduction, alternating with imports when renewable output is low.\\

\paragraph{France} demonstrates consistent public support for nuclear energy, with 75\% favoring continued operation of existing plants and $65\%$ supporting new construction according to a 2022 IFOP survey \cite{IFOP_Nuclear_Perception_2023}. Given the long-standing consensus on nuclear power as a strategic energy pillar and the country's established industrial base, significant policy reversals appear unlikely in the near term. The country electric grid relies heavily on nuclear power, which accounted for 67\% of its electricity generation in 2024 \cite{IEA_Electricity_2024}. Despite the substantial demand, which averaged $49.05 \pm 9.83$ GW in 2024, France is a net electricity exporter, with 21\% of total domestic production crossing the border in 2024.

\subsection{Selection of Reactor Models among existing Installations}
This study focuses on a curated set of reactor designs at an advanced stage of development.
The selection process begins with the the IAEA "Nuclear Reactors in the World" \cite{international2024iaea}, which catalogs operational, under-construction, and decommissioned reactors globally. \\

The initial selection process applied two primary exclusion criteria. First, designs were eliminated based on supplier origin, removing all reactors associated with countries unlikely to engage in nuclear technology transfer with Europe under the current geopolitical climate. This filter specifically excluded reactors developed in or supplied by Russia, China, and other nations maintaining restricted technology-sharing agreements with Western markets.

Second, reactors commissioned before 2000 were excluded from consideration, as their designs represent outdated technological standards. The complete filtering methodology is documented in the accompanying analysis notebook (PRIS Data Analysis.ipynb), which details how the initial dataset was reduced to nine distinct reactor models.

From this pre-selected group, three designs - the EPR-1750, APR-1400, and AP-1000 - were prioritized for comprehensive analysis. It should be noted that while the EPR (1650 MWe) and EPR-1750 share significant design similarities, they are treated as separate models in this study. Additionally, the Advanced Boiling Water Reactor (ABWR), though listed as operational in the IAEA ARIS database, was excluded following verification through the PRIS database, which indicates all ABWR units have suspended operation since 2011 \cite{IAEA_PRIS_2025}. The selection process and its outcomes are summarized in Table \ref{tab:reactor-selection}.


\begin{table}[!hb]
\centering
\begin{tabularx}{\textwidth}{llXXXX}
\toprule
\textbf{Reactor Name} & \textbf{Country} & \textbf{Technology-sharing} & \textbf{Commercial Operation} & \textbf{Up to Date} & \textbf{Decision} \\
\midrule
APR1400 & S.Korea & \cmark & \cmark & \cmark & \textbf{Selected} \\

CAREM Prototyp & Argentina & \cmark & \xmark & \cmark & Rejected \\

PRE KONVOI & Germany & \cmark & \cmark & \xmark & Rejected \\

EPR 1750 & France & \cmark & \cmark & \cmark & \textbf{Selected} \\

AP1000 & USA & \cmark & \cmark & \cmark & \textbf{Selected} \\

ABWR & Japan & \cmark & \xmark & \cmark & Rejected \\

APWR & Japan & \cmark & \xmark & \cmark & Rejected \\

EPR & France & \cmark & \cmark & \xmark & Rejected \\

OPR-1000 & S.Korea & \cmark & \cmark & \xmark & Rejected \\
\bottomrule
\end{tabularx}
\caption{Selection of reactor models for analysis}
\label{tab:reactor-selection}
\end{table}

\subsection{Selection of Reactors among emerging technologies}
Subsequently, the IAEA Advanced Reactors Information System (ARIS) database \cite{iaea2025aris}, was consulted to identify advanced reactor models classified as under design, under regulatory review, under construction, or operational. Demonstrations units and prototypes were excluded from this set, yielding 10 candidate designs. To avoid redundancy, models already included in the prior selection (e.g., Generation III+ designs captured in both datasets) were removed. Each remaining design was then assessed for development status, technological readiness, and applicability to European energy systems as summarized in Table \ref{tab:advanced-reactor-selection}.

\begin{table}[!hb]
\centering
\begin{tabularx}{\textwidth}{lllX}
\toprule
\textbf{Reactor Name} & \textbf{Country} & \textbf{Used} & \textbf{Reasoning} \\
\midrule

APR+ & S.Korea & \xmark & As reported in the manufacturer website \cite{kepcoenc2025} this design is mean to be an "indigenous reactor cleared of obstacles to export"\\

ATMEA & Japan, France & \xmark & Joint venture was abandoned in 2019 \cite{PLACEHOLDER}\\

BWRX300 & USA, Canada & \cmark & Design from GE-Hitachi, a lot of support from US and Canada as well as interes in the EU \cite{wnn2025bwrx300} \cite{wnn2025hungarysmr}.\\

i-SMR & S.Korea & \xmark & Never expressed interest in development outside south Asia and the middle east, no updates since early 2024.\\

iMSR400 & US & \xmark & No interest overseas, project seems to be at an early stage of development.\\

Natrium & US & \xmark & Demo plant \cite{terrapower2025natrium}.\\

Nuward & France & \xmark & After the change in direction of 2024 the project has been set back \cite{PLACEHOLDER}.\\

Nuscale & US & \cmark & Even if there was some financial controversy in 2024 \cite{PLACEHOLDER}. the project seems to be still moving forward in terms of goverment founding \cite{PLACEHOLDER} and regulatory review \cite{PLACEHOLDER}.\\

PRISM & US & \xmark & GE-Hitachi considers it a "concept ... currently being put into practice in two reactors: Natrium and ARC-100" \cite{PLACEHOLDER}.\\

Rolls Royce SMR & UK & \cmark & It's moving forward and from the technological point of view this is the simples as it it litteraly a scaled down PWR\cite{PLACEHOLDER}.\\

SMART & S.Korea, Saudi Arabia & \xmark & Plan to export only to south-east asia and middle east \cite{wnn2025smrdesign}.\\ \bottomrule

\end{tabularx}
\caption{Selection of advanced reactor models for analysis}
\label{tab:advanced-reactor-selection}
\end{table}

\subsection{Overview of the case study}
So we are going to consider 6 reactors to be deployed in 4 european countries....
The overall selection produced a set of six reactors, three Generation III+ and three SMRs. Their main characteristics are summarized in Table \ref{tab:final-reactor-selection}. The most notable differences are the integrated designs of the BWRX-300 and Nuscale SMR which eliminate the external primary loop in favour of steam generator inside the reactor pressure vessel. This allows for compact construction and allows to exploit natural circulation for cooling. Then the Nuscale SMR is the only model designed to be deployed in a multi unit configuration, with up to 12 modules in the same containment building, for this study the 6 module plant (VOYGR-6) has been selected as it recently recived the Standard Design Approval by the U.S. Nuclear Regulatory Commission in May 2025. "A standard design approval indicates that a proposed reactor design meets applicable agency safety requirements. Companies that seek to use the US460 design would have to file applications seeking permission to build and operate a nuclear reactor using the approved design" \cite{NRC_2025_033}.

\textcolor{red}{Maybe here add a paragraph for each design to explain their main features.}
\begin{table}[!hb]
\centering
\begin{tabularx}{\textwidth}{lclcX}
\toprule
\textbf{Reactor Name} & \textbf{Type} & \textbf{Construction} & \textbf{P$_e$ (MWe)} & \textbf{Supplier} \\
\midrule
APR1400 & PWR  & 1 Unit & 1340 & KEPCO (S.Korea) \\
EPR 1750 & PWR  & 1 Unit & 1630 & Framatome (France) \\
AP1000 & PWR  & 1 Unit & 1120 & Westinghouse (USA) \\
BWRX300 & i-BWR & 1 Unit & 300 & GE-Hitachi (USA-Canada)\\
Rolls Royce SMR & PWR  & 4 to 12 Units & 470 & Rolls Royce (UK) \\
Nuscale & i-PWR & 1 Unit & 77 & Nuscale Power (USA) \\
\bottomrule
\end{tabularx}
\caption{Final selection of reactor models for analysis}
\label{tab:final-reactor-selection}
\end{table}

